\section{Dataset Exploration and Cleaning}
\textbf{Class Distribution} \\
We balanced the training set across the six source
classes (about 4948 images each). This represents about a $1{:}5$ real-to-AI
ratio after merging (Fig.~\ref{fig:aspect} left). We keep this ratio and address
it during modelling through a class-weighted loss and a calibrated threshold. We
prefer this over re-sampling, because the evaluation splits have the same
imbalance.

\textbf{The important characteristic is aspect ratio} \\
Inspecting image sizes
reveals a close-to-perfect shortcut in train: real (COCO) images are nearly
always non-square (e.g.\ $640{\times}480$ or $640{\times}427$), whereas nearly
all AI images are square ($320{\times}320$ or $270{\times}270$). Amongst the
kept real images, $\RealNonSquare$ are non-square and only $\RealSquare$ are
square, while all $\AiSquare$ AI images are in fact square
(Fig.~\ref{fig:aspect} right). In contrast, the calibration, validation and
predict splits are uniformly $320{\times}320$. Aspect ratio is therefore an
artifact of the training split alone: a model that relies only on this metric
would label every square validation image as AI (possibly a 100\% false-positive
rate on real images). We remove this shortcut by resizing every training image
to a square, matching the test distribution --- the most important cleaning
decision for this dataset.\\
\textbf{Other observations}\\
The averaged FFT power spectrum
(Fig.~\ref{fig:spectrum}) indicates that real images consistently carry more
high-frequency energy than AI images. This is a genuine forensic cue that is
also present on validation, so it transfers. Secondly, AI images are slightly
brighter and less saturated. Finally, the augmented splits have smaller encoded
sizes (Fig.~\ref{fig:bytes}), indicating the augmentation pipeline applies
stronger JPEG compression. As shown in Task~1.3, this mostly destroys the
high-frequency cue and motivates augmentation-based robustness.

\begin{table}[h]
\centering\small
\begin{tabular}{lcc}
\toprule
Descriptive statistic (training) & Real (0) & AI (1--5)\\
\midrule
images & 4\,887 & 24\,489\\
typical resolution & $640{\times}480$ non-square & $320^2$ and $270^2$ square\\
mean encoded size & 50\,KB & 26\,KB\\
mean brightness (RGB) & 0.43 & 0.47\\
mean saturation & 0.29 & 0.28\\
\bottomrule
\end{tabular}
\caption{Per-class descriptive statistics of the training split (counts after
de-duplication). Resolution and encoded size are strongly correlated with the
label in train, but are neutralised by the square resizing; colour differences
are weak.}
\label{tab:desc}
\end{table}

\textbf{Deterministic cleaning pipeline (\texttt{clean.py})}\\
 We decoded all
training images and dropped undecodable streams ($\Nbad$ found), dropped exact
duplicates via an MD5 hash of the encoded bytes while keeping the first
occurrence ($\Ndup$ removed), and recorded per-image metadata. The result is a
deterministic cleaned index of $\Nkept$ rows. Squaring to the model resolution,
justified above, is applied in \texttt{prepare.py}. We intentionally avoid any
non-deterministic cleaning, or cleaning that alters image content (such as
denoising), which could erase the forensic cues the detector should use.

\begin{figure}[t]
  \centering
  \includegraphics[width=0.84\linewidth]{fig_class_aspect.png}
  \caption{Left: balanced source-class counts in train. Right: in train, real
  images are non-square and AI images are square --- a shortcut absent from the
  $320^2$ evaluation splits.}
  \label{fig:aspect}
\end{figure}

\begin{figure}[t]
  \centering
  \begin{subfigure}{0.47\linewidth}
    \includegraphics[width=\linewidth]{spectrum_validation.png}
    \caption{Radial power spectrum (validation).}
    \label{fig:spectrum}
  \end{subfigure}\hfill
  \begin{subfigure}{0.39\linewidth}
    \includegraphics[width=\linewidth]{fig_bytes.png}
    \caption{Encoded size by split.}
    \label{fig:bytes}
  \end{subfigure}
  \caption{Forensic cues. (a)~Real images keep more high-frequency energy in the
  radial power spectrum on validation; (b)~augmentation compresses images more
  heavily, as shown by the encoded size per split.}
\end{figure}
